% ═══════════════════════════════════════════════════════════════════════════
% MUSTERLÖSUNG: Examen Unidad 1 + 2 (Klasse 9) - Version 3.1
% ═══════════════════════════════════════════════════════════════════════════

\documentclass[a4paper,11pt]{article}

\usepackage[utf8]{inputenc}
\usepackage[T1]{fontenc}
\usepackage[ngerman]{babel}
\usepackage{helvet}
\renewcommand{\familydefault}{\sfdefault}

\usepackage[margin=18mm]{geometry}
\usepackage{xcolor}
\usepackage{tikz}
\usepackage{enumitem}
\usepackage{array}
\usepackage{booktabs}
\usepackage{multicol}
\usepackage{tabularx}
\usepackage{tcolorbox}
\usepackage{amssymb}
\usepackage{fancyhdr}
\usepackage{lastpage}

% Fachfarbe Spanisch
\definecolor{spanisch}{HTML}{c41e3a}
\definecolor{hellgrau}{HTML}{f5f5f5}
\definecolor{mittelgrau}{HTML}{888888}
\definecolor{loesung}{HTML}{c41e3a}
\definecolor{wahlfarbe}{HTML}{2c5aa0}

% Lösung
\newcommand{\lsg}[1]{{\color{loesung}\textbf{#1}}}

% Punktebox (rechts)
\newcommand{\punkte}[1]{\hfill\textcolor{mittelgrau}{\small /\,#1}}

% Aufgaben-Umgebung
\newcounter{aufgabe}
\newenvironment{aufgabe}[2]{%
    \refstepcounter{aufgabe}
    \vspace{4mm}
    \noindent\textbf{\color{spanisch}Ejercicio \theaufgabe: #1 \punkte{#2}}
    \vspace{1.5mm}
    \par
}{%
    \vspace{2mm}
}

\setlength{\parindent}{0pt}
\setlength{\parskip}{0pt}

% Seitennummerierung
\pagestyle{fancy}
\fancyhf{}
\renewcommand{\headrulewidth}{0pt}
\fancyfoot[C]{\small\color{mittelgrau}Seite \thepage\ von \pageref{LastPage}}

\begin{document}

% ═══════════════════════════════════════════════════════════════════════════
% KOPFBEREICH
% ═══════════════════════════════════════════════════════════════════════════

\noindent
\begin{tabularx}{\textwidth}{Xr}
    {\Large\bfseries\color{spanisch}Examen: Unidad 1 + 2 -- MUSTERLÖSUNG} & {\small Klasse 9 -- Spanisch} \\
\end{tabularx}

\vspace{1mm}
\noindent\rule{\textwidth}{0.4pt}
\vspace{3mm}

% ═══════════════════════════════════════════════════════════════════════════
% EJERCICIO 1: Vocabulario im Kontext
% ═══════════════════════════════════════════════════════════════════════════

\begin{aufgabe}{Vocabulario en contexto}{6}

\begin{enumerate}[itemsep=2mm]
    \item Du fragst nach \lsg{la cuenta} (die Rechnung).
    \item Du brauchst \lsg{la cuchara} (der Löffel).
    \item Die Milch stellst du in \lsg{la nevera} (den Kühlschrank).
    \item Du hängst \lsg{el cuadro} (das Bild) an die Wand.
    \item Du schaust in \lsg{el espejo} (den Spiegel).
    \item Bücher ordnest du in \lsg{la estantería} (das Regal).
\end{enumerate}

{\small\color{mittelgrau}Je 1 Punkt. Artikel muss korrekt sein!}

\end{aufgabe}

% ═══════════════════════════════════════════════════════════════════════════
% EJERCICIO 2: Fehlerkorrektur pedir/servir
% ═══════════════════════════════════════════════════════════════════════════

\begin{aufgabe}{Corrige los errores -- pedir y servir}{5}

\begin{enumerate}[itemsep=2mm]
    \item *pedo → \lsg{pido} (yo-Form: e→i im Stamm)
    \item *serven → \lsg{sirven} (ellos-Form: e→i im Stamm)
    \item *pedís tú → \lsg{pides} (tú-Form, nicht vosotros!)
    \item *servemos → \lsg{servimos} (nosotros: KEIN Stammwechsel!)
    \item *pede → \lsg{pide} (ella-Form: e→i im Stamm)
\end{enumerate}

{\small\color{mittelgrau}Je 1 Punkt. Erklärung nicht erforderlich, aber hilfreich für Teilpunkte.}

\end{aufgabe}

% ═══════════════════════════════════════════════════════════════════════════
% EJERCICIO 3: Conocer + a personal
% ═══════════════════════════════════════════════════════════════════════════

\begin{aufgabe}{El verbo conocer -- ¿con o sin \textit{a}?}{4}

\begin{enumerate}[itemsep=2mm]
    \item Yo \lsg{conozco a} tu hermano. (Person → mit a)
    \item ¿Tú \lsg{conoces --} Barcelona bien? (Stadt → ohne a)
    \item Nosotros \lsg{conocemos a} la profesora nueva. (Person → mit a)
    \item Ellos \lsg{conocen --} muchos restaurantes. (Orte → ohne a)
\end{enumerate}

{\small\color{mittelgrau}Je 1 Punkt (0,5 für Verb + 0,5 für a/--). Fehlendes \textbf{a} bei Person = 0,5P Abzug.}

\end{aufgabe}

% ═══════════════════════════════════════════════════════════════════════════
% EJERCICIO 4: Pretérito perfecto - gemischt
% ═══════════════════════════════════════════════════════════════════════════

\begin{aufgabe}{El pretérito perfecto}{6}

\begin{enumerate}[itemsep=2mm]
    \item Esta mañana yo \lsg{he trabajado} mucho. {\small(regulär: -ado)}
    \item ¿Qué \lsg{has hecho} (tú) hoy en clase? {\small(irregulär: hacer → hecho)}
    \item Mis padres nunca \lsg{han viajado} a Alemania. {\small(regulär: -ado)}
    \item Yo todavía no \lsg{he visto} la película. {\small(irregulär: ver → visto)}
    \item Nosotros \lsg{hemos escrito} una carta a la abuela. {\small(irregulär: escribir → escrito)}
    \item Hoy María no \lsg{ha comido} nada. {\small(regulär: -ido)}
\end{enumerate}

{\small\color{mittelgrau}Je 1 Punkt (0,5 haber + 0,5 Partizip).}

\end{aufgabe}

\newpage

% ═══════════════════════════════════════════════════════════════════════════
% EJERCICIO 5: Direktobjektpronomen im Dialog
% ═══════════════════════════════════════════════════════════════════════════

\begin{aufgabe}{Los pronombres de objeto directo}{5}

\begin{tcolorbox}[colback=hellgrau, colframe=mittelgrau, boxrule=0.5pt, arc=0pt]
\textbf{Mamá:} ¿Dónde está la lámpara nueva?\\
\textbf{Pablo:} \lsg{La} he puesto en el dormitorio. {\small(la lámpara = fem. Sg.)}\\[1mm]
\textbf{Mamá:} ¿Y los libros de la escuela?\\
\textbf{Pablo:} \lsg{Los} he dejado en la estantería. {\small(los libros = mask. Pl.)}\\[1mm]
\textbf{Mamá:} ¿Has visto mis gafas?\\
\textbf{Pablo:} Sí, \lsg{las} he visto en la cocina. {\small(las gafas = fem. Pl.)}\\[1mm]
\textbf{Mamá:} ¿Y el sofá nuevo?\\
\textbf{Pablo:} Sí, papá y yo \lsg{lo} hemos subido esta mañana. {\small(el sofá = mask. Sg.)}\\[1mm]
\textbf{Mamá:} ¿Y las plantas del balcón?\\
\textbf{Pablo:} \lsg{Las} he regado antes de salir. {\small(las plantas = fem. Pl.)}
\end{tcolorbox}

{\small\color{mittelgrau}Je 1 Punkt pro korrektem Pronomen.}

\end{aufgabe}

% ═══════════════════════════════════════════════════════════════════════════
% EJERCICIO 6: Bewegungsverben + Präpositionen
% ═══════════════════════════════════════════════════════════════════════════

\begin{aufgabe}{Los verbos de movimiento -- elige la preposición}{5}

\begin{tabular}{|p{6cm}|c|p{4cm}|}
\hline
\textbf{Satz} & \textbf{Lösung} & \textbf{Erklärung} \\
\hline
Bajo las cajas \rule{0.8cm}{0.4pt} sótano. & \lsg{al} & hin zu (a+el) \\
\hline
Llevo el televisor \rule{0.8cm}{0.4pt} dormitorio. & \lsg{al} & hin zu (a+el) \\
\hline
Saco la leche \rule{0.8cm}{0.4pt} la nevera. & \lsg{de} & weg von (de+la) \\
\hline
Pongo los libros \rule{0.8cm}{0.4pt} la estantería. & \lsg{en} & Position (in) \\
\hline
Subo las maletas \rule{0.8cm}{0.4pt} piso. & \lsg{al} & hin zu (a+el) \\
\hline
\end{tabular}

\vspace{2mm}
{\small\color{mittelgrau}Je 1 Punkt. bajar/subir + a = Zielort; sacar + de = Herkunft; poner + en = Position.}

\end{aufgabe}

% ═══════════════════════════════════════════════════════════════════════════
% EJERCICIO 7: Leseverstehen
% ═══════════════════════════════════════════════════════════════════════════

\begin{aufgabe}{Comprensión de lectura}{8}

\textbf{Erwartete Antworten (auf Spanisch):}

\vspace{3mm}

\begin{enumerate}[itemsep=4mm]
    \item ¿Por qué han ido al restaurante? (1P)

    \lsg{Han ido para celebrar el cumpleaños de la madre.}

    {\small\color{mittelgrau}Auch akzeptiert: ``Para el cumpleaños de su madre.''}

    \item ¿Qué ha pedido cada persona de la familia? (3P)

    \lsg{El padre: una pizza con jamón. La madre: pasta con mariscos. El/La narrador/a: una ensalada grande.}

    {\small\color{mittelgrau}Je 1P pro Person. Nur ``pizza/pasta/ensalada'' = 0,5P.}

    \item ¿Por qué ha pedido el/la narrador/a una ensalada? (1P)

    \lsg{Porque no ha tenido mucha hambre.}

    {\small\color{mittelgrau}Auch: ``No tenía hambre.'' / ``Porque no tenía hambre.''}

    \item ¿Cuánto ha costado la cena en total (con propina)? (1P)

    \lsg{Cincuenta euros. / 50 euros.} {\small(45 + 5 = 50)}

    {\small\color{mittelgrau}Rechnung muss stimmen. Nur ``45 euros'' = 0P.}

    \item ¿Crees que la familia va a volver? ¿Por qué? (2P)

    \lsg{Sí, (creo que) van a volver porque la madre ha dicho que ha sido la mejor cena. / Porque el restaurante ha sido perfecto. / Porque todo ha estado delicioso.}

    {\small\color{mittelgrau}1P für Ja/Nein mit Begründung aus Text. 2P für ausführliche Begründung mit Textstelle.}
\end{enumerate}

\end{aufgabe}

\newpage

% ═══════════════════════════════════════════════════════════════════════════
% EJERCICIO 8: Schreibaufgabe (Wahlaufgabe)
% ═══════════════════════════════════════════════════════════════════════════

\begin{aufgabe}{Producción escrita}{10}

\textbf{Bewertungsraster (gilt für beide Varianten):}

\vspace{2mm}

\begin{tabular}{|l|c|p{9cm}|}
\hline
\textbf{Kriterium} & \textbf{Max.} & \textbf{Beschreibung} \\
\hline
Inhalt & 3 & Alle geforderten Elemente vorhanden und sinnvoll umgesetzt \\
\hline
Grammatik & 3 & Korrekte Verbformen, Artikel, Adjektivkongruenz \\
\hline
Kohäsion/Kohärenz & 2 & Logischer Aufbau, Verknüpfungen, roter Faden \\
\hline
Rechtschreibung & 2 & Korrekte Orthographie, Akzente \\
\hline
\end{tabular}

\vspace{4mm}

\textbf{Variante A -- Beispieldialog:}

\vspace{2mm}

{\small
\begin{tabular}{ll}
\textbf{Cam.:} & \lsg{Buenas tardes. ¿Qué desean?} \\
\textbf{Cli.:} & \lsg{Buenas tardes. De primer plato quiero una sopa, por favor.} \\
\textbf{Cam.:} & \lsg{Muy bien. ¿Y de segundo plato?} \\
\textbf{Cli.:} & \lsg{De segundo, pido una paella.} \\
\textbf{Cam.:} & \lsg{Perfecto. ¿Y para beber?} \\
\textbf{Cli.:} & \lsg{Una coca-cola, por favor. ¿Qué postres tienen?} \\
\textbf{Cam.:} & \lsg{Tenemos flan, helado y tiramisú.} \\
\textbf{Cli.:} & \lsg{De postre, quiero un helado de chocolate.} \\
\textbf{Cam.:} & \lsg{¿Algo más?} \\
\textbf{Cli.:} & \lsg{No, gracias. ¿Me trae la cuenta, por favor?} \\
\end{tabular}
}

\vspace{4mm}

\textbf{Variante B -- Anforderungen:}
\begin{itemize}[itemsep=0mm]
    \item Mind. 4 Räume mit Möbeln beschrieben
    \item Korrekter Vokabular (Räume + Möbel aus U2)
    \item 80+ Wörter erreicht
\end{itemize}

\vspace{4mm}

{\small\color{mittelgrau}\textbf{Abzüge:}
\begin{itemize}[itemsep=0mm, topsep=1mm]
    \item Weniger als 10 Redebeiträge / 80 Wörter: -1P Inhalt
    \item Schwere Grammatikfehler: -0,5P pro Fehler (max. -3P)
    \item Fehlende Kohäsion: -1P
    \item Orthographiefehler: -0,25P pro Fehler (max. -2P)
\end{itemize}}

\end{aufgabe}

\vspace{5mm}

\noindent\rule{\textwidth}{0.4pt}

\vspace{2mm}

\noindent
\begin{tcolorbox}[colback=hellgrau, colframe=spanisch, boxrule=1pt, arc=0pt, title={\color{white}\textbf{Notenschlüssel (49 Punkte)}}]
\begin{tabular}{lll|lll}
Note 1: & 47--49 P & (96\%) & Note 4: & 20--29 P & (40\%) \\
Note 2: & 40--46 P & (80\%) & Note 5: & 10--19 P & (20\%) \\
Note 3: & 30--39 P & (60\%) & Note 6: & 0--9 P & (<20\%) \\
\end{tabular}
\end{tcolorbox}

\vspace{3mm}

\noindent
\begin{tcolorbox}[colback=wahlfarbe!10, colframe=wahlfarbe, boxrule=1pt, arc=0pt, title={\color{white}\textbf{AFB-Verteilung}}]
\begin{tabular}{lll}
\textbf{AFB I} (Reproduktion): & Ej. 4 (teilw.) & $\approx$ 10\% \\
\textbf{AFB II} (Anwendung): & Ej. 1, 3, 4, 5, 6, 7 (teilw.) & $\approx$ 52\% \\
\textbf{AFB III} (Transfer): & Ej. 2, 7.5, 8 & $\approx$ 38\% \\
\end{tabular}
\end{tcolorbox}

\vfill

\noindent\rule{\textwidth}{0.4pt}
{\small\color{mittelgrau} Musterlösung -- Nur für Lehrkräfte | Klassenarbeit Unidad 1 + 2 | Klasse 9 | 49 Punkte}

\end{document}
